\documentclass{exam}
\usepackage[utf8]{inputenc}
\usepackage{graphicx}
\usepackage{dblfloatfix}
\makeatletter
\def\@seccntformat#1{%
  \expandafter\ifx\csname c@#1\endcsname\c@section\else
  \csname the#1\endcsname\quad
  \fi}
\makeatother
\pointname{punten}
\thispagestyle{empty}
\pagestyle{empty}


\begin{document}
\begin{titlepage}
\center
\includegraphics[scale=0.08]{astrid}
\linebreak
\linebreak
\linebreak
\vspace{1em}{\huge \bfseries  De Astrid Quiz }
\linebreak

\vspace{1em}{\LARGE  Door cultuur home astrid} 
            \linebreak
\vspace{1em}{\large 18 maart astrid}
	\begin{center}
	\fbox{\fbox{\parbox{5.5in}{\centering
	De Regels
	\begin{enumerate}
	\item{Zoals bij elke quiz zijn er geen hulpmiddelen toegestaan! Gelieve uw gsm op stil te zetten en in het zakje op de tafel te steken.}
	\item{De punten die bij elke vraag te verdienen zijn, zijn altijd aangegeven op het antwoorden blad.}
	\item{Er is een bierprijs voorzien. De bar zal het aantal online bijhouden. Je kunt dit dus ook meevolgen op het computerscherm bij de bar.}
	\end{enumerate}
	}}}
	\end{center}
\end{titlepage}

\newpage

\begin{center}
\fbox{\fbox{\parbox{5.5in}{\centering
Schrijf hier uw antwoorden op. Er ook zijn notitie bladeren meegegeven.}}}
\end{center}
\vspace{5mm}
 
\makebox[\textwidth]{Teamnaam:\enspace\hrulefill}
 
\vspace{5mm}
 
\makebox[\textwidth]{Teamnummer:\enspace\hrulefill}
\section{Ronde 1:}
\begin{questions}

\question[1] \enspace\hrulefill
\question[1] \enspace\hrulefill
\question[1] \enspace\hrulefill
\question[1] \enspace\hrulefill
\question[1] \enspace\hrulefill
\question[1] \enspace\hrulefill
\question[1] \enspace\hrulefill
\question[1] \enspace\hrulefill
\question[1] \enspace\hrulefill
\question[1] \enspace\hrulefill

\end{questions}
\begin{table}[!b]
\centering
\begin{tabular}{|l|l|l|l|l|l|l|l|l|l|l|}
\hline
Vraag       & 1 & 2 & 3 & 4 & 5 & 6 & 7 & 8 & 9 & 10 \\ \hline
max. punten & 1 & 1 & 1 & 1 & 1 & 1 & 1 & 1 & 1 & 1  \\ \hline
score       &   &   &   &   &   &   &   &   &   &    \\ \hline
\end{tabular}
\end{table}
\newpage


\section{Ronde 2: Nostalgie}
\begin{questions}

\question[2] 2003
\vspace{5mm}
\question[2] Lion king-1994; Hercules-1997; Mulan-1998; Fantasia 2000- 1999
\vspace{5mm}
\question[2] Courage The Cowardly Dog \bf{of} Courage, het bange hondje
\vspace{5mm}
\question[2] Free Souffriau
\vspace{5mm}
\question[2] Lizzie Mcquire
\vspace{5mm}
\question[2] Marsupilami
\vspace{5mm}
\question[2] Kinderen voor Kinderen (KvK)
\vspace{5mm}
\question[2] Mojo Jojo
\vspace{5mm}
\question[2] 61 afleveringen, \bf{maar we aanvaarden 59-63}, 20+20+21 = 61
\vspace{5mm}
\question[2] Haai

\end{questions}
\newpage


\section{Ronde 3: fetisches}
\begin{questions}

\question[2] Oude voornaam: Bouwdewijn, niewe voornaam: Bo ,achternaam: Vanspilbeek (alles moet fonetisch kloppen)
\vspace{5mm}
\question[2] De seksuele identiteit komt overeen met het biologisch geslacht waarmee zij geboren zijn. ()
\vspace{5mm}
\question[2] Klittenband
\vspace{5mm}
\question[2] Lesbian/gay/bisexueel/transgender/queer/intersex/asexueel/ al de rest
\vspace{5mm}
\question[2] 21/3=7 (EXACT!!!)
\vspace{5mm}
\question[2] Asexuele
\vspace{5mm}
\question[2] Burundi
\vspace{5mm}
\question[2] Ballonnen
\vspace{5mm}
\question[2] Het verlangen naar personen van alle genderidentiteiten en biologische geslachten
\vspace{5mm}
\question[2] Kyra Gantios

\end{questions}
\newpage


\section{Ronde 4: Reclame}
\begin{questions}

\question[2] met GAMMA,GAMMA
\vspace{5mm}
\question[2] Douwe Egberts
\vspace{5mm}
\question[2] Disneyland
\vspace{5mm}
\question[2] Jagermaester
\vspace{5mm}
\question[2] Vandenborre
\vspace{5mm}
\question[2] Hans anders
\vspace{5mm}
\question[2] Evian
\vspace{5mm}
\question[2] De lijn
\vspace{5mm}
\question[2] Pornhub
\vspace{5mm}
\question[2] Heineken

\end{questions}

\newpage


\section{Ronde 5: sport}
\begin{questions}

\question[2] Frédérik Deburghgraeve
\vspace{5mm}
\question[2] 3de
\vspace{5mm}
\question[2] Basketbal
\vspace{5mm}
\question[2] Blauw, geel, zwart, groen, rood
\vspace{5mm}
\question[2] 't kuipke
\vspace{5mm}
\question[2] Serve
\vspace{5mm}
\question[2] Golf
\vspace{5mm}
\question[2] Ultimate frisbee
\vspace{5mm}
\question[2] Biathlon
\vspace{5mm}
\question[2] 0-0

\end{questions}

\newpage


\section{Ronde 6: wereldreis}
\begin{questions}

\question[2] Gent-Sint pieters Station
\vspace{5mm}
\question[2] Val Thorens
\vspace{5mm}
\question[2] Dom van keulen
\vspace{5mm}
\question[2] Kopenhagen, De kleine zeemeermin
\vspace{5mm}
\question[2] christ church's great hall, England
\vspace{5mm}
\question[2] Forst van Brest
\vspace{5mm}
\question[2] Tracks of Hanoi/ doorway railroad/doorstep railroad
\vspace{5mm}
\question[2] Krater van Dwerze/ The door to hell / doorway to hell
\vspace{5mm}
\question[2] Rainbow mountain/Regenboogberg / Vinicunca
\vspace{5mm}
\question[2] Jeanneke Pis in brussel

\end{questions}
\newpage


\section{Ronde 7: bierlogo's}
\begin{questions}

\question[2] \enspace\hrulefill
\vspace{5mm}
\question[2] \enspace\hrulefill
\vspace{5mm}
\question[2] \enspace\hrulefill
\vspace{5mm}
\question[2] \enspace\hrulefill
\vspace{5mm}
\question[2] \enspace\hrulefill
\vspace{5mm}
\question[2] \enspace\hrulefill
\vspace{5mm}
\question[2] \enspace\hrulefill
\vspace{5mm}
\question[2] \enspace\hrulefill
\vspace{5mm}
\question[2] \enspace\hrulefill
\vspace{5mm}
\question[2] \enspace\hrulefill

\end{questions}
\begin{table}[!b]
\centering
\begin{tabular}{|l|l|l|l|l|l|l|l|l|l|l|}
\hline
Vraag       & 1 & 2 & 3 & 4 & 5 & 6 & 7 & 8 & 9 & 10 \\ \hline
max. punten & 2 & 2 & 2 & 2 & 2 & 2 & 2 & 2 & 2 & 2  \\ \hline
score       &   &   &   &   &   &   &   &   &   &    \\ \hline
\end{tabular}
\end{table}
\newpage






%ronde 7
\begin{center}
\fbox{\fbox{\parbox{5.5in}{\centering
Schrijf hier uw antwoorden op. Er ook zijn notitie bladeren meegegeven.}}}
\end{center}
\vspace{5mm}
 
\makebox[\textwidth]{Teamnaam:\enspace\hrulefill}
 
\vspace{5mm}
 
\makebox[\textwidth]{Teamnummer:\enspace\hrulefill}
\section{Ronde 8:}
\begin{questions}
\question[1] {\begin{flushleft}
{\large Voorbeeldvraag}
\end{flushleft}}

\end{questions}
\begin{table}[!b]
\centering
\begin{tabular}{|l|l|l|l|l|l|l|l|l|l|l|}
\hline
Vraag       & 1 & 2 & 3 & 4 & 5 & 6 & 7 & 8 & 9 & 10 \\ \hline
max. punten & 1 & 1 & 1 & 1 & 1 & 1 & 1 & 1 & 1 & 1  \\ \hline
score       &   &   &   &   &   &   &   &   &   &    \\ \hline
\end{tabular}
\end{table}
\newpage

%tusenronde 3
\begin{center}
\fbox{\fbox{\parbox{5.5in}{\centering
Schrijf hier uw antwoorden op. Er ook zijn notitie bladeren meegegeven.}}}
\end{center}
 
\vspace{5mm}
 
\makebox[\textwidth]{Teamnaam:\enspace\hrulefill}
 
\vspace{5mm}
 
\makebox[\textwidth]{Teamnummer:\enspace\hrulefill}
\section{Tafel ronde 2: Wie is het?}
Freddy heeft een paar mooie foto's gevonden van het home Astrid presidium. Kan jij raden wie het is?
s
\begin{questions}
\question[2] {\begin{flushleft}
{\large 1: Twenty-one}
\end{flushleft} }

\end{questions}

\begin{table}[!b]
\centering
\begin{tabular}{|l|l|l|l|l|l|l|l|l|l|l|}
\hline
Vraag       & 1 & 2 & 3 & 4 & 5 & 6 & 7 & 8 & 9 \\ \hline
max. punten & 2 & 1 & 1 & 1 & 1 & 1 & 1 & 1 & 1 \\ \hline
score       &   &   &   &   &   &   &   &   &   \\ \hline
\end{tabular}
\end{table}
\newpage
\section{Tafel ronde 3: raadsel ronde}
\begin{questions}
\question[4] {\begin{flushleft}
{\large Raadsel 1: Twenty-one}
\end{flushleft}
\begin{flushleft}
Dit raadsel komt van de thuisstad van Tibo (cultuur home Astrid). In Koksijde staan er veel wetgevingen op hoogte bouw, dus is er maar één echt super hoog gebouw: de Twenty-one. Rara-ra, het gebouw heeft 21 verdiepen. Nu woont de beste vriend van Freddy op de 21ste verdieping. Elke dag als het mooi weer is neemt hij de lift naar beneden en de trap naar boven. Echter wanneer het slecht weer is neemt hij de lift naar beneden en ook de lift naar boven.
\linebreak
De vraag is, hoe komt dit? Waarom neemt hij niet altijd de lift?
\linebreak
\end{flushleft}
}
\it Hint: op slechte dagen neemt hij een paraplu mee!
\fillwithlines{\stretch{1}}
%raadsel 2
\question[2]{\begin{flushleft}
{\large Raadsel 2: Vogel}
\end{flushleft} Teken een vogel door 2 rechte strepen te plaatsen.\center
\includegraphics[scale=1]{kip}}
\newpage
\question[2] {\begin{flushleft}
{\large Raadsel 3: Symbolen}
\end{flushleft}
Bepaal de volgende symbolen (mintens 2 aanvullen ,maar maximum 5).
\center
\includegraphics[scale=1]{symbolen}
}


\newpage
\question[2] {\begin{flushleft}
{\large Raadsel 4: Ontcijfer}
\end{flushleft}
\begin{flushleft}
Ontcijfer deze boodschap!
\end{flushleft}
\begin{center}
Oleacahoydnebp xef 3
\end{center}
\it Hint! Kijk naar het nummer van het raadsel!
\fillwithlines{3em}
}

\question[2]{\begin{flushleft}
{\large Raadsel 5: converteer het symbool naar een decimaal getal.}
\end{flushleft}
\begin{center}
\includegraphics[scale=0.2]{76}
\end{center}
\fillwithlines{2em}
}

\question[2]{\begin{flushleft}
{\large Raadsel 6: Freddy probeert 8 konoginnnen op een schaakbord te zetten, zonder dat ze elkaar in de weg staan. Een koningin kan elke kant uit, dus ze kan horizontaal, verticaal en diagonaal bewegen. }
\end{flushleft}
\begin{center}

\includegraphics[scale=0.03]{schaakbord2}
\end{center}
\newpage
\question[4]{
Er zijn 4 verschillende dozen gegeven, telkens met bijhorend opschrift:

\vspace{0.2cm}

Doos A: \framebox[0.4\textwidth][l]{Het punt zit niet in doos B}

\vspace{0.2cm}

Doos B: \framebox[0.4\textwidth][l]{Het punt zit niet in doos C}

\vspace{0.2cm}

Doos C: \framebox[0.4\textwidth][l]{Het punt zit in doos D}

\vspace{0.2cm}

Doos D: \framebox[0.4\textwidth][l]{Het punt zit in doos A}
\vspace{0.2cm}

Je mag maar 1 doos kiezen en openen!

\textbf{Er is een wiskundige manier voor deze vraag! men kan dit formaliseren met proposities en eventueel verzamelingenleer.}

{\em Opmerking:} om beide punten te krijgen moet je ook proposities geven!}
}
\newpage
De meer lugubere raadsels:
\question[2]{
Een man ligt dood op het platteland met een rugzak naast hem.

Wat is er gebeurd?
\fillwithlines{\stretch{1}}
}
\question[2]{
Er rijd een zwarte auto op een zwarte baan, de gebouwen zijn zwart en de bewoners hebben hun zwarte gordijnen gesloten. Er geen enkel licht, dus de autolichten en straatverlichting staan uit. De auto stopt en laat de zwarte kat oversteken. Hoe komt het dat de autobestuurder de zwarte kat zag?

Wat is er gebeurd?
\fillwithlines{\stretch{1}}
}
\end{questions}
\begin{table}[!b]
\centering
\begin{tabular}{|l|l|l|l|l|l|l|l|l|l|l|}
\hline
Vraag       & 1 & 2 & 3 & 4 & 5 & 6 & 7 & 8 & 9 \\ \hline
max. punten & 4 & 2 & 2 & 2 & 2 & 2 & 4 & 2 & 2 \\ \hline
score       &   &   &   &   &   &   &   &   &   \\ \hline
\end{tabular}
\end{table}

\end{document}